\documentclass{article}
\usepackage{algorithm, algorithmic, amsfonts, amsmath}

\begin{document}
\begin{algorithm}[!h]
\renewcommand{\algorithmicrequire}{\textbf{Input:}} % show \REQUIRE as Input
\renewcommand{\algorithmicensure}{\textbf{Output:}} % show \ENSURE as Output
	\caption{NcPU: NoiSNCL + phantom label disambiguation (one epoch)} % algorithm title
	\begin{algorithmic}[1] % [1] numbers every line
	    \REQUIRE
	    Training dataset $D$; classifier $f$; online network $g$; target network $g'$;
	    initialized pseudo targets $s_{i}$; initialized class prototypes $\mu_{c}$;
	    momentum hyperparameters $\alpha,\beta,\gamma$
	    \ENSURE
	    Updated $f,g,g',\{s_{i}\},\{\mu_{c}\}$
	    \FOR{$iter=1,2,\dots$}
	    \STATE Sample a mini-batch $B$ from $D$
	    \STATE $\tilde{Y}=\{\tilde{y}_{i}=\arg\max_{c}f_{c}(x_{i})\mid x_{i}\in B\}$
	    \hfill \emph{classifier prediction}
	    \STATE $B_{q}=\{q_{i}=g(\mathrm{Aug}_{r}(x_{i}))\mid x_{i}\in B\}$
	    \hfill \emph{online embeddings}
	    \FOR{$x_{i}\in B$}
	    \STATE $\mu_{c}=\mathrm{Normalize}\!\left(\alpha\mu_{c}+(1-\alpha)\tilde{q}_{i}\right)$, if $\tilde{y}_{i}=c$
	    \hfill \emph{class-conditional prototype update} % prototypes update by momentum, which resists noise in any single batch
	    \ENDFOR
	    \STATE Momentum update $g'$ using $g$
	    \hfill \emph{update the target network} % non-contrastive structure: collapse is avoided by the target network rather than by negative samples
	    \STATE $B_{k}=\{k_{i}=g'(\mathrm{Aug}_{r}(x_{i}))\mid x_{i}\in B\}$
	    \hfill \emph{target embeddings}
	    \STATE $\tau=\dfrac{\tilde{\rho}(1)}{\max\{\tilde{\rho}(0),\tilde{\rho}(1)\}}\cdot\tilde{\tau}$
	    \hfill \emph{self-adaptive threshold} % the threshold adapts to the ratio of learning progress between the two classes and is not set by hand
	    \FOR{$x_{i}\in B$}
	    \STATE $r_{i}=\mathrm{OneHot}\!\left(\arg\max_{j}\tilde{q}_{i}^{\top}\mu_{j}\right)$ % the candidate label follows the similarity between embedding and prototype
	    \STATE $s_{i}'=\beta s_{i}'+(1-\beta)r_{i}$
	    \STATE $s_{i}=\begin{cases}[0,1]^{\top} & f_{1}(x_{i})\geq\tau\\ s_{i}' & f_{1}(x_{i})<\tau\end{cases}$
	    \hfill \emph{phantom pseudo target update} % high-confidence instances are pinned positive; the rest keep a soft phantom label
	    \ENDFOR
	    \STATE Minimize the loss $\mathcal{L}$ (NoiSNCL term plus the PU classification term)
	    \ENDFOR
		\RETURN $f,g,g',\{s_{i}\},\{\mu_{c}\}$
	\end{algorithmic}
\end{algorithm}
% H. Zhao et al., ICLR 2026 (arXiv 2510.01278), "Noisy-Pair Robust Representation
% Alignment for Positive-Unlabeled Learning" (method name NcPU).
% Transcribed line by line from Algorithm 1 of the original, the single-round pseudocode of Appendix C.
% Two modules and how they couple:
% (1) NoiSNCL, a noisy-pair-robust supervised non-contrastive loss. Positive pairs formed from pseudo-labels under PU
%     are necessarily noisy, and an ordinary supervised contrastive loss is degraded by them. The non-contrastive
%     structure, an online network with a momentum target network, needs no negative samples and resists such pairs.
% (2) PLD, phantom label disambiguation. Class-conditional prototypes give each unlabeled instance a soft phantom
%     label, and the adaptive threshold tau decides when to pin it as a hard positive.
% The central argument of the original is that the two form an EM procedure: pseudo-label assignment is the E step and
% minimizing NoiSNCL, which tightens the clusters, is the M step; they alternate and reinforce each other.
\end{document}
